\subsection{Archive layout and integration}

The archive contains a complete source snapshot needed by the supplied
tests, an additive patch against the audited commit, the article sources
and rendered PDF, raw measurements, certificate-bearing audit records,
compact source cases, reproduction scripts, and provenance notes. The
\path{fast/} snapshot preserves the pinned files and adds the new modules.
Its sibling \path{reports/} directory contains only historical oracle
dependencies required by the existing tests.

The patch uses paths rooted at \path{Topology/UnknotRecognition/fast/}.
It adds the feature modules, tests, and research drivers. It does not alter
the production recognition schedule. The new article and the retained
data are separate reviewable artifacts; the repository's intake process
can place them under its next report number.

From a checkout at the pinned revision, inspect and check the patch:
\begin{lstlisting}
git apply --stat /path/to/integration.patch
git apply --check /path/to/integration.patch
\end{lstlisting}
The package build verifies that the patch applies to the pinned source
and reproduces its included added files. A later repository revision
should be checked independently for path conflicts before applying it.

\subsection{Python setup and tests}

The new native algorithms use the Python standard library. The maintained
project supports Python 3.10 or later. This research environment uses
Python 3.12.14; use Python 3.12 for the pinned optional package set.
Regina is optional for most native tests but required to
repeat the advertised independent Regina comparisons. The archive gives
the exact optional package versions in \path{requirements-research.txt};
it does not bundle third-party binary distributions.

A reproducible environment can be prepared at the extracted archive root:
\begin{lstlisting}
python3 -m venv .venv
. .venv/bin/activate
python -m pip install -r requirements-research.txt
python reproduce/run_checks.py --full
\end{lstlisting}
Without the optional dependencies, native feature checks can be run with:
\begin{lstlisting}
python reproduce/run_checks.py
\end{lstlisting}
The latter command uses standard-library test discovery and verifies
retained native proofs. Its output identifies optional comparisons
that were skipped. For the explicit repository regression:
\begin{lstlisting}
cd fast
python -m unittest discover -s tests -v
\end{lstlisting}

\subsection{Minimal source-bound example}

From the supplied \path{fast/} directory:
\begin{lstlisting}[language=Python]
from fastunknot import Diagram
from fastunknot.normal_transport import transport_seed_decide
from fastunknot.normal_transport_verify import (
    verify_transport_disk_certificate,
)

diagram = Diagram.from_braid(2, [1])
answer = transport_seed_decide(
    diagram, strategy="score", max_work=2_000_000
)
assert answer["status"] == "UNKNOT"
assert verify_transport_disk_certificate(
    diagram, answer["certificate"]
)
\end{lstlisting}

The function accepts optional boundary shellings and initial minimum-span
optimization. It performs no subsequent cohomology solve after a move.
A failed or resource-limited search returns \code{INCONCLUSIVE}. The
certificate schema is \code{diagram-transport-disc-v1}; it includes the
source PD, canonical exterior, optional shelling transcript, initial
heights, every triangulation and cochain-transport proof, final coordinates,
and a positive essential-disc component certificate.

No rank-one or primitivity claim about the initial heights is needed by
the positive consumer: the independently established essential disc in
the verified exterior is sufficient. The producer uses a primitive class
to make its search useful. This separation keeps the consumer's
mathematical responsibilities smaller.

\subsection{Low-level geometric and boundary APIs}

\begin{lstlisting}[language=Python]
from fastunknot.cocycle_transport import (
    bipyramid_cocycle_score,
    cocycle_collapse_candidates,
    descend_cocycle,
)
score = bipyramid_cocycle_score([0, -1, 1, 3, 4])
assert score["euler_loss"] == 4
assert score["normal_disc_increase"] == 5
\end{lstlisting}

The five-height input order is belt first, then apices. A transport proof
uses schema \path{pachner-cocycle-transport-v1}. The mathematical formulas
are valid for arbitrary integral heights; the native geometry wrapper
has the compact orientable one-torus-boundary contract stated in
Section~\ref{sec:cochains}.

\begin{lstlisting}[language=Python]
from normal_orbit_research.fixtures import layered_torus
from fastunknot.marked_boundary import normal_marked_boundary_order
from fastunknot.marked_boundary_verify import (
    verify_normal_marked_boundary_certificate,
)

triangulation, coordinates = layered_torus(8)
marks = [0]
answer = normal_marked_boundary_order(
    triangulation, coordinates, marks,
    weight_encoding="moments", record_certificate=True
)
assert answer["status"] == "COMPLETE"
assert verify_normal_marked_boundary_certificate(
    triangulation, coordinates, marks, answer["certificate"]
)
\end{lstlisting}

The lower-level \code{marked_boundary_order(size, pairings, marks, ...)}
accepts a validated raw degree-two multigraph encoding. The reference
\code{weight_encoding="one_hot"} is retained. The default moment mode
changes the query encoding, not the returned cyclic solution.
An optional \code{start_half_edge} fixes the direction of one marked
component. Unmarked circles remain present in the output histogram.

\subsection{Recreate audits and benchmarks}

From \path{fast/}, the following commands write new records into
\path{../reproduced/}; they preserve the supplied measurements:
\begin{lstlisting}
python -m transport_research.audit local \
  --output ../reproduced/cocycle-local.json
python -m transport_research.audit corpus \
  --corpus ../reproduce/cocycle-source-cases.json \
  --output ../reproduced/cocycle-corpus.json
python -m transport_research.audit benchmark \
  --sizes 16 32 64 128 256 --rounds 5 \
  --output ../reproduced/cocycle-benchmark.json
python -m transport_research.splitting \
  --output ../reproduced/cocycle-splitting.json
python -m marked_boundary_research.matched_final \
  --output ../reproduced/marked-boundary-paired.json
\end{lstlisting}
The local transport audit uses Regina unless \code{--no-regina} is
specified. The final paired boundary driver records individual samples
and exact equality checks. The broader native-boundary experiment can
also be repeated with:
\begin{lstlisting}
python -m marked_boundary_research.benchmark \
  --output ../reproduced/marked-boundary.json
\end{lstlisting}
Timing benchmarks should run sequentially with other heavy jobs paused.
Exact numerical outputs, counts, and accepted certificates are
reproducibility targets; machine-dependent elapsed times are not expected
to match bit for bit.

\subsection{Build the article}

The included table and figure sources reproduce the paper's displays from
the retained JSON measurements. The already generated figures permit a
PDF rebuild without Python plotting packages. From \path{article/}:
\begin{lstlisting}
latexmk -pdf -interaction=nonstopmode -halt-on-error \
  unknot_geometric_transport.tex
\end{lstlisting}
Rebuilding the tables and plots first uses:
\begin{lstlisting}
python reproduce/build_tables.py
\end{lstlisting}
from the archive root; the plotting dependency is listed separately in the
research environment file. The final PDF was rendered to page images and
visually inspected for broken equations, table overflow, and pagination.
